\section{Native cohomology of the actual normalized cyclic cover}

Let $X$ be integral and let $g:R\to X$ be the original raw cyclic
cover attached to the original invertible line $L$, degree $n>0$
and section $s$.  Let
$\phi:R_\eta\to R$ be its actual base-generic-fibre projection.
The normalized cover used here is the native relative normalization
$Z=\phi^{\mathrm{nor}}$, with the literal map
$\pi=\phi.\mathrm{fromNormalization}\circ g:Z\to X$.
The cohomology statements below concern its structure sheaf and
actual pulled-back lines.

\begin{lemma}[Affineness of the native normalized projection]
\label{lem:kn-normalized-projection-affine}
\leanok
The composite of the native relative-normalization projection with
an affine morphism is affine.  In particular the actual $\pi$ is
affine for the original invertible $L$ and positive degree $n$.
\end{lemma}
\begin{proof}
The native normalization projection is integral, hence affine.
Affine morphisms are closed under composition.  For the actual
cyclic cover the original raw projection is finite by its original
local monic degree-$n$ equations, so it is affine.  The actual
generic-fibre map is quasi-compact and quasi-separated by base
change from the original generic-point map; these are precisely
the witnesses defining its native normalization.  Applying the
composite statement therefore gives the asserted actual map.
\end{proof}

\begin{lemma}[Quasicoherence of the actual structure pushforward]
\label{lem:kn-normalized-structure-quasicoherent}
\leanok
\uses{lem:kn-normalized-projection-affine}
The native module $\pi_*\mathcal O_Z$ is quasicoherent.
\end{lemma}
\begin{proof}
The actual structure module on any scheme is invertible, hence
quasicoherent.  Pushforward along the actual affine $\pi$ preserves
quasicoherence, proving the assertion.
\end{proof}

\begin{lemma}[Native positive cohomology across the actual affine projection]
\label{lem:kn-normalized-affine-cohomology}
\lean{EnvelopeProperty.AffinePushforwardCohomology.subsingleton_H_iff_pushforward, EnvelopeProperty.AffinePushforwardCohomology.subsingleton_H_twist_iff_pushforward}
\leanok
\uses{lem:kn-normalized-projection-affine}
Suppose the original base is separated and has a finite affine open
cover.  For a quasicoherent module $M$ on $Z$, positive native sheaf
cohomology vanishes if and only if the same positive native
cohomology of $\pi_*M$ vanishes.  For an invertible $T$ on $X$, the
same statement identifies $M\otimes\pi^*T$ with $\pi_*M\otimes T$.
In particular
\[
 H^q(Z,\pi^*T)=0
 \quad\Longleftrightarrow\quad
 H^q(X,\pi_*\mathcal O_Z\otimes T)=0\qquad(q>0).
\]
On a compact separated base the finite affine cover is obtained
internally.
\end{lemma}
\begin{proof}
The inverse images of the original affine opens and their finite
intersections are affine because $\pi$ is affine.  Quasicoherence
therefore makes positive cohomology vanish on these opens.  The
native sheaf-cohomology/ordered-Cech comparison on the original
cover identifies both positive groups with the same ordered Cech
homology, since pushforward has exactly the original sections on
these inverse images.  This proves the first equivalence.  The
actual line projection formula identifies the twisted pushforward
with $\pi_*M\otimes T$, giving the second equivalence.  The
structure module is invertible and thus quasicoherent; cancel its
native tensor-unit isomorphism to obtain the pulled-line statement.
Compactness gives finitely many affine opens covering the base;
choose a finite ordering and apply the same proof.
\end{proof}

For the weighted statements, let $X$ be smooth of relative
dimension $d$ over a perfect field $k$, with $n$ a unit in $k$,
and use an original SNC presentation of the zero ideal of $s$.
These are hypotheses on the original base and section.

\begin{definition}[The original normalized coefficient family]
\label{def:kn-normalized-coefficients}

Give the actual zero ideal $J$ of the same original section an
original SNC presentation $J=\prod_j D_j^{a_j}$.  This presentation
itself gives the whole SNC witness.  For $i\in\{0,\ldots,n-1\}$,
take the original native coefficient module
$E_i=L^{-i}\otimes\mathcal Hom(\mathcal O_X(-\lfloor(i/n)\operatorname{div}(s)\rfloor),
\mathcal O_X)$, using that same zero ideal and intrinsic rounding.
\end{definition}

\begin{lemma}[Whole weighted cohomology and the actual summand retracts]
\label{lem:kn-normalized-weighted-cohomology}
\leanok
\uses{def:kn-normalized-coefficients}
The existing native normalized eigensheaf isomorphism from PR~215
is an isomorphism of the whole original pushforward
\[
 \pi_*\mathcal O_Z\simeq\bigoplus_{i=0}^{n-1}E_i,
 \qquad
 E_i=L^{-i}\otimes
 \mathcal Hom\bigl(\mathcal O_X(-\lfloor(i/n)\operatorname{div}(s)\rfloor),
                   \mathcal O_X\bigr).
\]
For any invertible $T$, this gives a whole twisted isomorphism with
$\bigoplus_i(E_i\otimes T)$.  Each actual summand is a retract, and
in every native cohomological degree $q$ the whole twisted
pushforward cohomology vanishes exactly when all the actual
$H^q(X,E_i\otimes T)$ vanish.
\end{lemma}
\begin{proof}
Apply the original tensor-by-$T$ functor to the existing actual
normalized weighted isomorphism.  Tensoring by an invertible line
is an equivalence, hence preserves the finite coproduct; this
constructs the stated whole twisted isomorphism.  In the finite
coproduct take the original inclusion of a chosen summand and the
projection which is its identity and zero on the other summands.
Their composite is the identity.  Conjugating these maps by the
actual whole isomorphism gives a retraction on the actual twisted
pushforward.  Applying the native underlying-sheaf and cohomology
functors preserves this retraction, so zero whole cohomology forces
zero cohomology of every original summand.  Conversely native
cohomology is additive: the identity on a finite biproduct is the
sum of its projection-inclusion composites.  If the cohomology of
each summand is zero, its image under cohomology makes every such
composite zero, hence makes the identity of the whole cohomology
group zero.  The coproduct-biproduct isomorphism gives the result
for the original finite coproduct.
\end{proof}

\begin{theorem}[The actual normalized cyclic positive-cohomology criterion]
\label{thm:kn-normalized-positive-cohomology}
\leanok
\uses{lem:kn-normalized-affine-cohomology,lem:kn-normalized-weighted-cohomology}
For the actual original normalized cover and every $q>0$,
\[
 H^q(Z,\pi^*T)=0
 \quad\Longleftrightarrow\quad
 \forall i<n,\ H^q(X,E_i\otimes T)=0.
\]
The finite affine cover may be supplied explicitly, or obtained
internally when the original separated base is compact.
\end{theorem}
\begin{proof}
Compose the actual affine pulled-line comparison with the whole
weighted-cohomology equivalence.  Every positive natural degree has
the form $q_0+1$, so these comparisons cover precisely all asserted
degrees.  On a compact base use the internally obtained finite
affine cover in the first comparison.
\end{proof}

The displayed equivalence preserves the unknown cover cohomology;
it supplies no vanishing theorem for that cover and no canonical
finite-duality identification.
